% !TeX spellcheck = de_DE

Eine zentrale Voraussetzung des Feedbackloops ist die automatisierte Bewertung der erzeugten Texte.
Hierfür werden verschiedene Metriken verwendet, die unterschiedliche Aspekte der Textqualität erfassen sollen.
Dazu zählen die in dem vorherigem Kapitel \ref{ch:metrikenanalyse} genannten Lesbarkeitsmetriken, sprachliche Strukturmerkmale sowie regelbasierte Anforderungen an Leichte Sprache Plus.

Da bisher keine fixe Methodik existiert um Leichte Sprache Texte sicher zu Bewerten, untersucht diese Arbeit zusätzlich die Eignung verschiedener Metriken für den beschriebenen Anwendungsfall.
Die Bewertungskomponente wird auch dafür verwendet, zu analysieren, welche Metriken sinnvoll genutzt werden können und welche Kennzahlen einen geeigneten Indikator für die Qualität von LS+-Texten darstellen.